\section{考察}
本章では，結果章で得られた学習曲線の傾向について，
(i) 方策選択（バンディット）手法の特性，
(ii) 報酬設計の変更による影響
の観点から考察する．

\subsection{方策選択手法による性能差}
報酬設計1において，PRQ系列はいずれのタスクでもBaselineを上回り，
方策再利用が学習初期の性能向上に寄与することが確認できた．
一方で，PRQ内の方策選択手法
（Boltzmann，$\varepsilon$-greedy，UCB）の優劣は
タスクに依存して変化した．
Task 5では$\varepsilon$-greedyが立ち上がりが速く高い性能を示したが，
エラーバーが比較的大きく，
試行間のばらつきが大きい傾向も見られた．
$\varepsilon$-greedyは探索が一様であるため，
偶然有効な方策を早期に引き当てた場合には急速に性能が伸びる一方，
引き当てのタイミングが遅れる試行では学習が遅れ，
分散が大きくなりやすいと解釈できる．

一方，Task 12やTask 15ではUCBおよびBoltzmannが優位であり，
$\varepsilon$-greedyは相対的に低い性能に留まった．
UCBは推定価値に不確実性（選択回数の少なさ）を加点することで，
未探索の方策を段階的に試す方向へ選択を誘導する．
そのため，タスクによって有効な方策が変化しやすい条件や，
初期の推定値が不安定な条件では，
$\varepsilon$-greedyよりも探索が安定し，
高い性能に繋がった可能性がある．
Boltzmannも確率的に複数方策を試しながら価値差に応じて
選択確率を変化させるため，
UCBほどの探索保証はないものの，
温度の設定次第で探索と活用のバランスを滑らかに調整でき，
タスクによってはUCBに近い挙動を示したと考えられる．

\subsection{報酬設計の変更による影響}
報酬設計2では，ステップ負報酬および壁衝突ペナルティを導入したため，
累積平均利得が負の領域から推移するタスクが多く見られた．
この変化は「ゴール到達だけでなく，
途中の行動コストを明示的に最適化対象へ含めた」ことに対応している．
しかし，利得の絶対値が変化しても，
PRQ系列がBaselineを一貫して上回る傾向は維持されており，
方策再利用の効果は報酬構造の変更に対して比較的頑健であることが
示唆される．

一方で，PRQ内の3手法の関係は報酬設計によって変化した．
例えばTask 5では報酬設計2でBoltzmannが最も高い利得を示し，
UCBが相対的に低くなった．
これは，負報酬を導入したことで短期的なコストの違いが価値推定に
反映されやすくなり，価値差に応じて確率的に選択を集中させる
Boltzmannが有利に働いた可能性がある．
一方，Task 9やTask 12ではUCBが高い性能を維持しており，
ペナルティ導入により探索が難しくなった状況でも，
不確実性を考慮した選択が有効なケースが存在することが分かる．
このことから，報酬設計は「どの方策選択が有利か」
を変えうる要因であり，単一条件で得た結論を一般化するには
注意が必要である．

\subsection{限界と今後の課題}
本研究の比較は，迷路環境および設定した2種類の報酬設計に基づくものであり，
報酬スケールや迷路構造，方策ライブラリの多様性が変化した場合に
同様の傾向が再現されるかは追加検証が必要である．
また，本研究ではグリッドサーチによりパラメータを選定したが，
実運用を想定するとタスクごとに広範な探索を行うことはコストが高い．
今後は，環境の特徴量や学習進行に応じてパラメータを自動調整する手法
（適応的温度制御や探索率制御）を導入し，
チューニング負荷を低減することが課題となる．
さらに，関数近似を用いる大規模状態空間や深層強化学習への拡張，
複数のソース環境からの方策選択なども検討対象となる．
